
% \subsection{Architektura i koncepcje Kubernetesa}
% 
% Kubernetes, będący otwartym systemem do automatyzacji wdrażania, skalowania oraz zarządzania aplikacjami
% kontenerowymi, stanowi obecnie dominującą platformę orkiestracyjną w ekosystemach chmurowych. Jego
% architektura opiera się na modelu \emph{master--node}, w którym \textbf{klaster Kubernetesa} składa się
% z \gls{controlplane}, zazwyczaj uruchamianej na węzłach typu \emph{master}, oraz z wielu \gls{workernode},
% na których działają kontenery z aplikacjami.
% 
% \subsubsection{Podstawowe komponenty: control plane, nodes, pods, deployments, services}
% 
% \gls{controlplane} obejmuje kluczowe komponenty systemowe, takie jak \textbf{kube-apiserver},
% który udostępnia interfejs API Kubernetesa i stanowi centralny punkt komunikacji w klastrze;
% \textbf{kube-scheduler}, odpowiedzialny za przydzielanie nowo tworzonych \gls{pod} do odpowiednich
% węzłów; \textbf{kube-controller-manager}, zarządzający zestawem kontrolerów realizujących logikę
% utrzymania pożądanego stanu systemu (m.in. kontroler replikacji czy kontroler punktów końcowych);
% oraz \textbf{etcd}, będący rozproszoną bazą danych typu klucz--wartość, przechowującą aktualny stan
% klastra.
% 
% Każdy \gls{workernode} zawiera komponent \textbf{kubelet}, pełniący rolę agenta komunikującego się
% z płaszczyzną sterowania oraz zarządzającego cyklem życia podów i kontenerów na danym węźle,
% a także \textbf{kube-proxy}, który odpowiada za implementację reguł sieciowych dla \gls{service},
% umożliwiając komunikację pomiędzy podami oraz dostęp do usług zewnętrznych.
% 
% \subsubsection{Mechanizmy rozszerzeń: CRDs i operatory}
% 
% Kubernetes operuje na zestawie abstrakcji, takich jak \gls{pod}, stanowiące najmniejszą jednostkę
% wdrożeniową i grupujące jeden lub więcej kontenerów; \gls{deployment}, zapewniające deklaratywne
% zarządzanie stanem aplikacji, w tym automatyzację aktualizacji oraz wycofywania zmian; oraz
% \gls{service}, będące abstrakcjami sieciowymi definiującymi logiczny zbiór podów oraz politykę
% dostępu do nich.
% 
% Elastyczność platformy jest dodatkowo rozszerzana poprzez mechanizmy takie jak \gls{crd}, które
% umożliwiają definiowanie własnych obiektów API i integrację niestandardowych zasobów z natywnym
% modelem sterowania Kubernetesa. Na ich bazie budowane są \gls{operator}, czyli komponenty
% oprogramowania enkapsulujące wiedzę domenową i automatyzujące zarządzanie złożonymi aplikacjami
% oraz ich cyklem życia w klastrze. Zrozumienie tych fundamentalnych elementów stanowi podstawę do
% analizy mechanizmów autonomicznego sterowania, które wykorzystują i rozszerzają natywne możliwości
% platformy Kubernetes.

\subsection{Architektura i koncepcje platformy Kubernetes}

Kubernetes (często zapisywany jako K8s) to platforma orkiestracji kontenerów typu open-source, 
służąca do automatyzacji wdrażania, skalowania i zarządzania aplikacjami. System ten wywodzi się 
bezpośrednio z wewnętrznego systemu zarządzania klastrami firmy Google o nazwie Borg 
\cite{verma2015large} i opiera się na dekadzie doświadczeń w uruchamianiu systemów produkcyjnych 
na masową skalę.

Fundamentalną cechą Kubernetesa jest jego deklaratywny model konfiguracji. W przeciwieństwie do 
podejścia imperatywnego, gdzie administrator definiuje sekwencję kroków do wykonania, w 
Kubernetesie definiuje się \textit{stan pożądany} (ang. \textit{desired state}) systemu. Platforma, 
poprzez ciągłe pętle kontrolne (ang. \textit{control loops}), autonomicznie podejmuje działania 
zmierzające do uzgodnienia stanu rzeczywistego ze zdefiniowanym \cite{burns2019kubernetes}.

\subsubsection{Komponenty płaszczyzny sterowania i węzłów roboczych}

Architektura klastra opiera się na modelu \textit{master-node}, w którym wyróżniamy warstwę sterującą 
(\gls{controlplane}) oraz warstwę wykonawczą (\gls{workernode}).

\textbf{\Gls{controlplane}} (płaszczyzna sterowania) zarządza stanem klastra i podejmuje globalne 
decyzje (np. o planowaniu zadań). W jej skład wchodzą:
\begin{itemize}
    \item \textbf{kube-apiserver}: Centralny komponent udostępniający interfejs REST API. Stanowi 
      "front-end" dla płaszczyzny sterowania i jako jedyny komunikuje się bezpośrednio z bazą 
      danych klastra \cite{kubernetes_docs}.
    \item \textbf{etcd}: Spójna, wysokodostępna baza danych typu klucz-wartość (ang. \textit{key-value store}), 
      wykorzystywana do trwałego przechowywania wszystkich danych klastra, w tym konfiguracji i stanu 
      obiektów.
    \item \textbf{kube-scheduler}: Komponent odpowiedzialny za obserwację nowo utworzonych podów, 
      które nie mają przypisanego węzła, i wybór optymalnego \gls{workernode} do ich uruchomienia. 
      Decyzja ta opiera się na analizie dostępnych zasobów, polityk \textit{affinity} oraz ograniczeń 
      sprzętowych \cite{k8s_scheduling}.
    \item \textbf{kube-controller-manager}: Proces uruchamiający podstawowe pętle kontrolne klastra 
      (np. Node Controller, Replication Controller). Każdy kontroler jest odpowiedzialny za obserwację 
      konkretnego aspektu stanu klastra i reagowanie na odchylenia.
\end{itemize}

\textbf{\Gls{workernode}} (węzeł roboczy) to maszyna (fizyczna lub wirtualna), na której uruchamiane są 
aplikacje. Kluczowe komponenty węzła to:
\begin{itemize}
    \item \textbf{kubelet}: Agent działający na każdym węźle, który przyjmuje specyfikacje \gls{pod} 
      (zwykle od API servera) i zapewnia, że opisane w nich kontenery są uruchomione i działają poprawnie. 
      Realizuje on lokalną pętlę uzgadniania stanu na poziomie węzła \cite{kubernetes_architecture}.
    \item \textbf{kube-proxy}: Komponent sieciowy obsługujący reguły sieciowe na węźle. Umożliwia 
      komunikację sieciową do podów z wnętrza i spoza klastra, implementując mechanizm \gls{service}.
    \item \textbf{Container Runtime}: Oprogramowanie odpowiedzialne za uruchamianie kontenerów 
      (np. containerd, CRI-O).
\end{itemize}

\subsubsection{Abstrakcje i obiekty API}

Kubernetes wprowadza warstwę abstrakcji nad infrastrukturą sprzętową, operując na obiektach logicznych:

\begin{itemize}
    \item \textbf{\Gls{pod}}: Najmniejsza, atomowa jednostka wdrożeniowa w Kubernetesie. Pod reprezentuje pojedynczą instancję procesu w klastrze i może składać się z jednego lub wielu ściśle powiązanych kontenerów, współdzielących zasoby pamięci masowej i sieciowej (wspólny adres IP).
    \item \textbf{\Gls{deployment}}: Obiekt wyższego poziomu zarządzający replikowanymi aplikacjami. Umożliwia deklaratywne aktualizacje (np. \textit{Rolling Update}) oraz skalowanie liczby replik Podów. To właśnie na poziomie Deploymentu najczęściej definiowane są limity zasobów (CPU/RAM), które są przedmiotem optymalizacji w niniejszej pracy.
    \item \textbf{\Gls{service}}: Abstrakcja definiująca logiczny zbiór Podów oraz politykę dostępu do nich (zazwyczaj jako stały adres IP i nazwa DNS). Zapewnia to luźne powiązanie (ang. \textit{loose coupling}) między zależnymi mikroserwisami.
\end{itemize}

\subsubsection{Mechanizmy rozszerzeń: CRD i Wzorzec Operatora}

Jedną z najistotniejszych cech architektury Kubernetes jest jej rozszerzalność. Mechanizm \textbf{\gls{crd}} pozwala użytkownikom na definiowanie nowych typów zasobów, które są obsługiwane przez API server na równi z zasobami natywnymi (takimi jak Pod czy Service) \cite{kubernetes_crd}.

Połączenie niestandardowego zasobu (\gls{crd}) z dedykowanym kontrolerem, który w pętli monitoruje ten zasób, dało początek tzw. \textbf{Wzorcowi Operatora} (ang. \textit{Operator Pattern}). W literaturze przedmiotu Operator definiowany jest jako metoda pakowania, wdrażania i zarządzania aplikacją Kubernetes, która programowo kodyfikuje wiedzę operacyjną niezbędną do utrzymania usługi \cite{dobies2020operators}.
W kontekście niniejszej pracy, koncepcja ta jest kluczowa – projektowany \textit{Smart Tuner} jest w istocie operatorem, który wykorzystuje \gls{crd} do przechowywania konfiguracji i autonomicznie steruje parametrami standardowych mechanizmów \gls{hpa} oraz limitami zasobów, zastępując w tym procesie człowieka.

